\section{Conclusion}
\label{sec:conclusion}

I study how the cost of family space shapes fertility and what housing
policy can do about it. I take the decline in fertility since 2007 as given,
as a fall in the taste for children, and ask how housing policy changes
births from 2023 on. Households choose whether to have children, whether and
how much to rent or own, and how much to save; buying requires a down
payment, earnings are risky, and developers enter freely to build.

The space a child needs is a recurring cost. It is paid for as long as the
child lives at home, and it costs about the same every period whether the
household rents or owns: at the 2007 price, about 2.9\% of mean annual
earnings a year. Because this cost is paid while earnings are uncertain,
earnings risk weighs on fertility. Lowering the standard deviation of
earnings shocks by about a quarter, with mean earnings unchanged, raises
completed fertility by 12\% at the 2007 price.

Policy raises births when it lowers the cost of a child rather than the cost
of a home, and its effects build up over a generation. Removing the down
payment for all households raises births by 1.3\% on impact, but young
households use the credit to buy the small homes they would otherwise rent,
and the cohort aged 18 to 21 in 2023 ends with 1.1\% fewer children. A
housing supply expansion with the same effect on impact keeps raising
births, to 3.0\% above the no-policy path by 2063 and 4.3\% in the long run:
every cohort faces cheaper space, and the additional children form
households of their own. The same credit restricted to households with a
child raises births by 5.1\% on impact and the completed fertility of the
2023 cohort by 4.9\%, because it raises the value of having a child rather
than the value of owning.

A higher property tax shows the same mechanism from the other side. It
lowers the house price by 7.0\% and helps young households buy, but it
raises the yearly cost of holding a room, and rents rise by 14.5\%. Births
barely move on impact and rise by 2.2\% in the long run, and this gain comes
from the rebate: returned equally to all households, the revenue moves
resources toward low- and middle-earning households, who have the additional
children. Returned in proportion to the rooms each household holds, the same
tax lowers births.

Two lessons follow. Access to ownership is not what keeps families out of the
space they need, so policies that make buying easier do little for births.
And because births today shape the number of households a generation later,
policies that look alike on impact can differ widely over time. Housing
policy aimed at fertility must therefore be judged by its long-run effects.
